\documentclass[11pt]{article}
\usepackage[margin=1in]{geometry}
\usepackage{amsmath,amssymb,amsthm,mathtools}
\usepackage{booktabs}
\usepackage{verbatim}
\usepackage{multicol}
\usepackage[colorlinks=true,linkcolor=blue,citecolor=blue,urlcolor=blue]{hyperref}
\hypersetup{pdftitle={201 Operations for Rank-48 4x4 Matrix Multiplication over Rings Containing 1/2},pdfauthor={Greg Sidebottom}}
\newtheorem{theorem}{Theorem}[section]
\newtheorem{lemma}[theorem]{Lemma}
\newtheorem{proposition}[theorem]{Proposition}
\theoremstyle{remark}
\newtheorem{remark}[theorem]{Remark}
\newcommand{\Q}{\mathbb{Q}}
\newcommand{\Z}{\mathbb{Z}}
\newcommand{\F}{\mathbb{F}}
\newcommand{\Zh}{\Z[\tfrac12]}
\newcommand{\Pt}{P^{\mathsf T}}
\newcommand{\ops}{\operatorname{ops}}
\title{201 Operations for Rank-48 $4\times4$ Matrix Multiplication\\
over Rings Containing $1/2$}
\author{Greg Sidebottom}
\date{Draft --- October 22, 2026}

\begin{document}
\maketitle

\begin{abstract}
Rank-48 bilinear schemes multiply $4\times4$ matrices with $48$
multiplications over any ring in which $2$ is invertible.  Applied
recursively, a scheme whose three linear maps use $q$ additions and
scalings costs $(1+q/32)\,n^{\log_4 48}+o(n^{\log_4 48})$ operations,
so $q$ fixes the leading constant.  The best published count was $204$
(Dumas, Pernet, Sedoglavic, Stapleton and Tichavsk\'y, 2026), lowered to
$202$ by R.~Patel's value-set search.  We give a scheme in the same
class whose maps use $55+57+89=201$ operations, leading constant
$7+9/32$, with no negations, so the count is the same whether or not
sign changes are charged.  It was found by annealing over the class's
symmetry orbit for representatives with a low operation floor and few
non-unit coefficients, then running a fast reimplementation of Patel's
large-neighbourhood search seeded from the matrices alone, with the
output map optimized through its transpose.  The scheme is checked by
PLinOpt's tools, by an independent exact verifier, by Patel's checker,
and by two Lean~4 proofs that it computes $AB$ over every commutative
ring in which $2$ is invertible.  No representative we tried goes below
$201$, and about twenty reach it exactly.  To test whether the programs
found so far can be recombined into a $200$, we pose that question
exactly, as the existence of a closed subset of the union of their
signal sets with one element fewer, and encode it for SAT solvers.  The
answer is no for all three maps, both at the record point and at a
second, balanced $201$ point, so no recombination of the programs found
so far reaches $200$.  The hardest question, the record's output map,
was settled by refuting a $125$-vector part of its $137$-vector
dictionary and then twelve cases, one for each remaining vector.
Encoding acyclicity by a total order instead of by depth levels made
these refutations up to $17$ times faster.
\end{abstract}

\section{Introduction}

Strassen's algorithm applied to $2\times2$ blocks of $2\times2$
matrices multiplies $4\times4$ matrices with $49$ multiplications over
any ring~\cite{Strassen69}.  In characteristic two the record is
$47$~\cite{Fawzi22,KM23}.  In characteristic zero it is $48$: AlphaEvolve
found a rank-48 scheme with complex coefficients~\cite{Novikov25}, and
Dumas, Pernet and Sedoglavic turned it into one with rational
coefficients, valid over any ring containing $1/2$~\cite{DPS25,DPS26a}.
The rational schemes known today all descend from AlphaEvolve's complex
one; Kaporin's independent complex rank-48 scheme~\cite{Kaporin24} has
no real equivalent~\cite{MSY26}.

Used recursively, a scheme's multiplication count fixes the exponent and
its additions fix the constant.  For an
$\langle n_0\times n_0\times n_0:r\rangle$ scheme whose linear maps use
$q$ additions and scalings in total,
\[
F(n) = r\,F(n/n_0) + q\,(n/n_0)^2
\quad\Longrightarrow\quad
F(n) = \Bigl(1+\frac{q}{r-n_0^2}\Bigr)n^{\log_{n_0} r} + o\bigl(n^{\log_{n_0} r}\bigr),
\]
so at $r=48$, $n_0=4$ the constant is $1+q/32$.  For rank $48$, $q$
fell from $315$~\cite{DPS25,PLinOpt} to $204$~\cite{DPSST26}, and to
$202$ in October 2026, when
Ruchit Patel contributed a $202$-operation program to
PLinOpt~\cite{PLinOpt} and published the search that found
it~\cite{Patel26}.

\paragraph{Contributions.}
\begin{itemize}
\item A rank-48 scheme over $\Zh$ whose linear maps use
  $55+57+89=201$ operations (constant $7+9/32\approx7.281$), with no
  negations; Appendix~\ref{app:scheme} lists its programs.  It was
  contributed to PLinOpt (data file \texttt{4x4x4\_48\_201}) and
  merged on October 22, 2026 (Section~\ref{sec:scheme}).
\item The method (Section~\ref{sec:found}): orbit annealing for a
  representative with a low operation floor and few non-unit
  coefficients; Patel's value-set large-neighbourhood search,
  reimplemented in Rust with PLinOpt's objective and seeded from the
  matrices alone; and a successive-halving sweep over the
  representatives.  Applied to a sparser representative, seeded from the
  matrix, and run many times over many representatives, Patel's method
  reaches $201$.
\item Five independent checks, including two Lean~4 proofs over every
  commutative ring with $2$ invertible (Section~\ref{sec:verify}).
\item Evidence that $201$ is a plateau for this class under value-set
  search, and an exact formulation of the question whether the programs
  found so far can be recombined into a cheaper one, with SAT
  encodings and verdicts (no recombination reaches $200$ at either
  point), and the work still in progress
  (Section~\ref{sec:lower}).
\end{itemize}

\section{Records}\label{sec:records}

Table~\ref{tab:rank48} lists the operation counts for rank-48 schemes
without change of basis.  With a change of basis the recursive cost
drops further, to constants between $5+2/3$ and $6+13/32$ depending on
the encoding~\cite{DPSST26}; that is a different cost model, and this
paper does not use it.

\begin{table}[ht]
\centering
\caption{Operations (additions and scalings) of rank-48 schemes for
$\langle4,4,4\rangle$ over rings containing $1/2$, without change of
basis, and the resulting constant of $n^{\log_4 48}$.}\label{tab:rank48}
\begin{tabular}{rlll}
\toprule
operations & constant & split $L+R+P$ & source \\
\midrule
$315$ & $10+27/32$ & & \cite{DPS25}, PLinOpt data~\cite{PLinOpt} \\
$204$ & $7+3/8$ & & \cite{DPSST26} \\
$202$ & $7+5/16$ & $56+56+90$ & R.~Patel, PLinOpt data~\cite{PLinOpt,Patel26} \\
$\mathbf{201}$ & $\mathbf{7+9/32}$ & $55+57+89$ & this paper, PLinOpt data~\cite{PLinOpt} \\
\bottomrule
\end{tabular}
\end{table}

\section{Model}\label{sec:model}

A rank-$r$ scheme for $\langle4,4,4\rangle$ is a triple of matrices
$L\in\Zh^{r\times16}$, $R\in\Zh^{r\times16}$, $P\in\Zh^{16\times r}$
with
\[
\operatorname{vec}(AB) = P\bigl((L\operatorname{vec}A)\odot(R\operatorname{vec}B)\bigr),
\]
where $\operatorname{vec}$ lists entries row by row and $\odot$ is the
entrywise product.  Product $k$ is
$m_k=\langle U_k,A\rangle\langle V_k,B\rangle$ with $U_k,V_k$ the
$4\times4$ matrices in row $k$ of $L$ and $R$, and $AB=\sum_k m_kW_k$
with $W_k$ in column $k$ of $P$.  The identity is equivalent to the
$16^3=4096$ Brent equations
$\sum_k P_{c,k}L_{k,a}R_{k,b}=[\,a,b\text{ contribute to }c\,]$, a
statement over $\Q$.  When they hold, the identity holds for matrices
over any ring in which $2$ is invertible, commutative or not, because
each product keeps its $A$-factor on the left; the scheme therefore
recurses on blocks.

\paragraph{Cost.}  We use PLinOpt's model~\cite{PLinOpt}.  A linear
straight-line program computes the rows of a matrix $M$ from its inputs
by lines $t\coloneqq u\pm v$ (one operation) and $t\coloneqq c\,u$ for a
constant $c\ne\pm1$ (one operation, a \emph{scaling}); copies and
negations are free.  $\ops(M)$ is the fewest operations of such a
program, and a scheme costs $\ops(L)+\ops(R)+\ops(P)$.  Patel's model
also charges each negation; the programs in this paper contain none, so
their counts agree in the two models.

\paragraph{Transposition.}  The output side is handled through its
transpose.  By the transposition principle~\cite{BCS97} a linear program
for a matrix without zero rows or columns reverses into one for its
transpose with the same scalings, and for $\Pt\in\Zh^{48\times16}$
\[
\ops(P)=\ops(\Pt)+(48-16)=\ops(\Pt)+32 .
\]
All three sides are then instances of one problem: compute $48$ vectors
of $\Zh^{16}$ from the $16$ unit vectors.

\paragraph{Floor.}  A row that is, up to sign, not a unit vector needs an
operation of its own, so $\ops(M)\ge|T(M)|$, where $T(M)$ is the set of
distinct such rows up to sign.  The \emph{floor} of a scheme is
$|T(L)|+|T(R)|+|T(\Pt)|+32$.

\paragraph{Symmetry.}  If $(U_k,V_k,W_k)_k$ is a scheme, so are
\[
(X^{\mathsf T}U_kY^{-\mathsf T},\;Y^{\mathsf T}V_kZ^{-\mathsf T},\;
X^{-1}W_kZ)_k
\qquad\text{and}\qquad
(a_kU_k,\;b_kV_k,\;(a_kb_k)^{-1}W_k)_k
\]
for invertible $X,Y,Z$ (substitute $A\mapsto XAY^{-1}$,
$B\mapsto YBZ^{-1}$)~\cite{deGroote78} and nonzero scalars $a_k,b_k$
(rescale product $k$).  If $X,Y,Z\in\mathrm{GL}_4(\Zh)$, that is, with
determinant $\pm2^e$, and each $a_k,b_k$ is $\pm$ a power of $2$, every
coefficient stays in $\Zh$.  All points of an orbit have
the same $48$ multiplications but different linear maps, and therefore
different floors and operation counts; choosing the point is half of the
optimization.

\section{The 201-operation scheme}\label{sec:scheme}

Table~\ref{tab:sides} summarizes the scheme.  Its coefficients are
small: $L$ and $R$ have entries in $\{0,\pm1,\pm\tfrac12\}$, with $32$
entries $\pm\frac12$ each, and $P$ has entries in $\{0,\pm1,\pm2\}$,
with $8$ entries $\pm2$; $72$ coefficients in all are not in
$\{0,\pm1\}$.  The floor is $156=42+44+38+32$, and the programs spend
$45$ operations above it, $13$ on each input side and $19$ on the
output side.  Each program line is one operation: a binary addition or
subtraction, or, four times in all, a halving.

\begin{table}[ht]
\centering
\caption{The 201-operation scheme, side by side.  \emph{Unit rows} are
rows that are a single coefficient $\pm1$ (free); \emph{targets} are the
remaining rows, all distinct; \emph{depth} is the longest chain of
dependent operations.}\label{tab:sides}
\begin{tabular}{lrrrrrr}
\toprule
side & unit rows & targets & additions & halvings & operations & depth \\
\midrule
$L$ & $6$ & $42$ & $53$ & $2$ & $55$ & $13$ \\
$R$ & $4$ & $44$ & $55$ & $2$ & $57$ & $13$ \\
$\Pt$ & $10$ & $38$ & $57$ & $0$ & $57$ & $21$ \\
$P$ & & & $89$ & $0$ & $89=57+32$ & $21$ \\
\midrule
total & & & $197$ & $4$ & $\mathbf{201}$ & \\
\bottomrule
\end{tabular}
\end{table}

\paragraph{Signs.}  Negating row $k$ of $L$ (or of $R$) together with
column $k$ of $P$ gives another valid scheme, so the sign of every
product's factors is free.  The programs as found contain output
negations, which PLinOpt does not charge but Patel's model does; we
re-emitted them with signs chosen so that no line negates anything.
The resulting triple, a point of the same orbit, costs $201$ in both
models, and it is the one in Appendix~\ref{app:scheme} and in PLinOpt.

\paragraph{Other 201s.}  The count is not tied to one point.  About
twenty annealed points reach $201$, in one of three splits:
$55+57+89$, its mirror image $57+55+89$, and $56+56+89$.  The last
appears at a balanced point with floor $150=40+40+38+32$ and $104$
non-unit coefficients, the profile of Patel's own representative
(Section~\ref{sec:lower}).

\section{How the scheme was found}\label{sec:found}

The search ran on October 20--21, 2026, on one Apple M4~Pro machine.
Table~\ref{tab:history} lists the milestones.

\subsection{Choosing a representative}

The annealer moves through the orbit by elementary steps on one of
$X,Y,Z$: add $c$ times one row to another with $c\in\{\pm1,\pm\frac12,
\pm2\}$, swap or negate rows, or scale a row by $2$ or $\frac12$, keeping
every coefficient in $\{0,\pm\frac14,\pm\frac12,\pm1,\pm2,\pm4\}$.  It
minimizes the floor plus $\lambda$ times the number of coefficients not in
$\{0,\pm1\}$ (Metropolis acceptance, restarts, a pool of the best
distinct points, each re-verified on all $4096$ Brent equations).  The
penalty matters: a low floor with many halves and doubles is expensive
to program.  The runs found three kinds of points: the profile of the
published $204$ representative (floor $146$ but $136$ non-unit
coefficients), that of Patel's $202$ representative (floor $150$, $104$),
and \emph{sparse} points (floor $156$--$159$, $68$--$72$).  The 201 comes
from the $\lambda=0.3$ pool of sparse points.

\subsection{Local refinement is not enough}

Our first pipeline refined PLinOpt's optimizer output at a fixed point
by re-solving small sub-programs.  It settled at $238$ on the best
sparse point and at $253$ on Patel's point, whose stored programs cost
$202$.  An exact solver for windows of up to $16$ consecutive program
lines, with up to four helper values, gained $3$ more ($235$), found
nothing on the other sides, and showed the other windows optimal among
programs with at most two helpers.  So most of the gap was not in how
a sub-program computes its values but in which values are computed at
all, which is what Patel's method searches over.

\subsection{Value-set search}

Patel's search~\cite{Patel26} identifies a program with its set of
signal vectors: the inputs, the targets, and the intermediate values,
each a vector of $\Zh^{16}$.  A set is \emph{closed} if its members can
be computed one at a time from the inputs, each by one operation on
members already computed ($u\pm v$, $2u$, $u/2$, or a rescaling of a
proportional member); a closed set yields a program with one operation
per non-input member, so the cost is the set's size minus $16$.  Each step of the search is a large-neighbourhood
move~\cite{Shaw98}: it proposes one to eight new vectors from six
families (residuals of targets against members, coordinated packages of
residuals, sparse pivots, residuals among members, halvings and
doublings, and broad random combinations), adds them, and then deletes
members in random order while the set stays closed.  Steps that lower
the cost are accepted, steps that keep it are accepted unless the new
set is in a tabu list of recent sets, and steps that raise it by one are
accepted with a small probability.

\subsection{A faster implementation, and seeds from the matrix}

Patel's Python spends most of a step in a sign planner for his cost
model, which charges negations; PLinOpt's model does not need it.  We
reimplemented the search in Rust (\texttt{src/bin/lns48.rs}): vectors in
fixed point over $1/64$, closure and deletion on bitsets, PLinOpt's
objective, and an exact re-verification of every program written.  It
runs about $80$ steps per second per thread, $70$ times the original.

The second change was the seed.  Every run seeded with programs derived
from PLinOpt's optimizer, Python or Rust, converged to $56$ on the
sparse point's $L$; seeded instead with the naive program that computes each
row of the matrix as a chain of additions (pruned by the same deletion
before the search starts), the search reached $55$ in twenty seconds.
The seed decides the basin, so we replaced long runs by many short
restarts from matrix-only seeds: twelve seeds of three minutes per side
reached $202=55+58+89$ on the sparse point.  The output side is always
searched as $\Pt$, a $48\times16$ map like $L$ and $R$, and transposed
back with PLinOpt's \texttt{transpozer}.

\subsection{The sweep}

The last step was a successive-halving sweep~\cite{JT16} over $90$
annealed points: $30$ seconds per side for every point, then $120$
seconds for the best $20$, then $600$ seconds with two seeds for the
best $6$.  Two points reached $201$ in the second stage, and all six
finalists ended at exactly $201$.

\begin{table}[ht]
\centering
\caption{Milestones of the search, October 20--21, 2026.  Counts are
PLinOpt's (additions and scalings).}\label{tab:history}
\small
\begin{tabular}{rlp{8.8cm}}
\toprule
total & $L+R+P$ & step \\
\midrule
$238$ & $70+67+101$ & sparse point, PLinOpt optimizer plus pair-level re-solving \\
$235$ & $67+67+101$ & exact window search on $L$ \\
$205$ & $56+58+91$ & Patel's search~\cite{Patel26} on the same point, $P$ through $\Pt$ \\
$204$ & $56+58+90$ & Rust reimplementation \\
$203$ & $55+58+90$ & seeded from the matrix alone \\
$202$ & $55+58+89$ & twelve short restarts per side (ties Patel's 202) \\
$\mathbf{201}$ & $55+57+89$ & sweep over $90$ annealed points \\
\bottomrule
\end{tabular}
\end{table}

\section{Verification}\label{sec:verify}

The triple in Appendix~\ref{app:scheme} passes five independent checks.
\begin{enumerate}
\item PLinOpt's \texttt{SLPchecker} accepts each program against its
  matrix and reports $53+2$, $55+2$ and $89+0$ additions plus scalings.
\item PLinOpt's \texttt{MMchecker} accepts the three matrices as a
  matrix multiplication scheme, in rational mode, on random inputs
  (five runs, three more with $40$-bit values).
\item Our verifier \texttt{matmul/dps48/slp48check.py} evaluates the
  programs exactly over $\Q$ and checks all $4096$ Brent equations
  ($0$ failures), counting operations in both cost models.
\item Patel's own checker~\cite{Patel26} reports $201$ linear
  operations, $4096$ tensor identities, and a well-formed program DAG.
\item Two Lean~4~\cite{Lean4} proofs with Mathlib~\cite{Mathlib}, checked
  by Lean's kernel.  \texttt{Matmul48\_201.correct}
  (\texttt{proof/Mm/Correct48\_201.lean}) takes the three programs line
  by line as hypotheses, $201$ operation wires and $48$ products with
  each halving written $u\cdot h$, and proves that the $16$ outputs are
  the entries of $AB$.  \texttt{rank48\_op201\_scheme\_eq\_mul}
  (\texttt{proof/Mm/Ring\_4x4\_48\_201.lean}) states the result in
  Mathlib's own terms:
\begin{verbatim}
theorem rank48_op201_scheme_eq_mul (h : R) (h2 : 2 * h = 1)
    (A B : Matrix (Fin 4) (Fin 4) R) :
    rank48_op201_scheme h A B = A * B
\end{verbatim}
  for every commutative ring \texttt{R}, where the scheme is a function
  of $201$ commented \texttt{let} bindings and $48$ products.  Each
  entry is closed by \texttt{linear\_combination}~$e\cdot(2h-1)$, with
  the polynomial $e$ computed by the generator through exact division by
  $2h-1$.  The proofs check in about $15$ and $90$ seconds and use only
  Lean's three standard axioms.
\end{enumerate}
Validity over non-commutative rings follows from check~3, since the
Brent equations are exactly the condition for the bilinear identity with
non-commuting entries.

\section{Can 201 be lowered? Work in progress}\label{sec:lower}

\subsection{A plateau}

No route we tried goes below $201$:
\begin{itemize}
\item about twenty distinct annealed points with near-minimal penalized
  floor reach exactly $201$, in one of the three splits above, and none
  goes lower; of the $19$ such points within two orbit moves of the
  record point, three reach $201$ and the rest $203$;
\item other profiles, from annealer variants that weight the three sides
  differently: points with side floors $42/42/42$ reach $L+R=110$ but
  $P=93$ ($203$), and random orbit walks that ignore the floor cost
  $212$ or more;
\item longer and wider searches at the record point: hot annealing
  ($330{,}000$ steps), proposal batches of $16$ and $24$ vectors with
  eight deletion orders, Metropolis acceptance at two temperatures,
  searching $P$ directly as a $16\times48$ map, and replacing the random
  deletion by an exact minimum closed subset (a mixed-integer program
  solved by HiGHS~\cite{HiGHS} over the set plus $60$ proposals, about
  $2000$ exact steps): $R$ and $\Pt$ stay at $57$ and $P$ at $89$
  throughout.
\end{itemize}
Two regularities hold at the floor-$150$ and floor-$156$ points:
$\ops(\Pt)=|T(\Pt)|+19$, so $P$ costs $|T(\Pt)|+51$ (also at the
$42/42/42$ points), and $\ops(L)+\ops(R)=112$.

\subsection{Recombination as an exact question}

Independent runs end in different programs with the same cost.  Could a
mixture of them do better?  Let the \emph{dictionary} $D$ be the union
of the signal sets of a collection of programs for one side, and let a
\emph{production} of $v\in D$ be one operation that computes $v$ from
other members, found by enumerating all pairs of members (sums,
differences, proportional rescalings).  The question
\begin{quote}
is there a closed set $S$ with inputs $\subseteq S\subseteq D$, all
targets in $S$, and at most $c$ non-input members?
\end{quote}
asks whether the side can be programmed in $c$ operations using only
vectors that some program in the collection computes.  Because the
productions of a subset are determined by its own members, a negative
answer for a dictionary is also a negative answer for every subset of
it.  The question can be enlarged by adding vectors to $D$ that no
program used: search proposals, or a \emph{structured} family, all
sub-vectors of at most three terms of each target (with the target's
coefficients) together with all pairwise sums and differences of
targets.

\subsection{Encodings}

The first formulation was a mixed-integer program with a selection
variable $x_v$ per vector, a variable $y_p$ per production, and
continuous levels with big-M constraints for acyclicity, solved by
HiGHS~\cite{HiGHS}.  Its feasibility form proves small cases in a
minute, but its relaxation is weak: four runs on dictionaries of
$86$--$137$ vectors reached their six-hour limits undecided.

The SAT encoding (\texttt{lns48 --merge ... --cap c --cnf}) uses the same
$x_v$ and $y_p$ with clauses
\[
x_t\ (t\text{ a target}),\qquad
y_p\to x_u\ (u\text{ an operand of }p),\qquad
x_v\to\textstyle\bigvee_{p\text{ produces }v}y_p,
\]
a sequential counter for $\sum x_v\le c$~\cite{Sinz05}, and three
strengthenings that are sound because any solution can be reduced to
one that satisfies them: $y_p\to x_v$, at most one production per
vector, and every selected non-target is an operand of a selected
production.  Acyclicity is encoded in one of two ways.
\begin{itemize}
\item \emph{Levels}: $g_{v,j}$ means that $v$ has depth at least $j$, for
  $j\le\ell$, with clauses that place each selected production's result
  above its operands.  This is complete only if $\ell$ is at least the
  depth of some optimal program; $\ell=c$ always is.
\item \emph{Order}: one variable per pair of non-input vectors and two
  transitivity clauses per triple encode a total order, and each
  selected production puts its operands before its result.  There is no
  depth bound.
\end{itemize}
A satisfying assignment is decoded into a program and re-verified.  We
validated the encoding on questions with known answers: on the record
point's $L$ dictionary it answers no at $54$ and yes at $55$, decoding to
a correct $55$-operation program; and the structured dictionary for $R$
at $c=57$ is satisfiable, with the record's own program as the model.
The solvers are CaDiCaL~3.0.0~\cite{CaDiCaL} and kissat~4.0.4~\cite{Kissat}.

\subsection{Verdicts}

Table~\ref{tab:verdicts} lists the verdicts.  At the record point no
recombination of the programs found lowers $L$ or $R$, for $R$ even with
$200$ added proposals; at the balanced point none lowers any of the
three sides, so $201=56+56+89$ is final there for this question.  The
record's $\Pt$ was the hard case.  Its programs are deep, with chains of
$21$ dependent operations against $13$ on $L$ and $R$, because rows of
$\Pt$ are built from one another in long chains.  The level encoding
pays for that depth.  Bounded depth is easy: no $56$-operation program of
depth at most $6$ exists ($0.09$ seconds), nor of depth at most $10$
($748$ seconds).  With the full depth the question stayed open for eight
hours under both solvers and every split we tried.  The order encoding
removes the depth variables altogether, and it refutes the balanced
point's $\Pt$ in $306$ seconds, against $88$ minutes with levels.

On the record's $\Pt$ dictionary, built from $32$ programs, the order
encoding refutes the parts contributed by the first $22$, $24$ and $28$
programs in $33$, $406$ and $5000$ seconds.  The time grows about
twelvefold per $14$ vectors, which puts the whole $137$-vector question
near twelve hours.  The refuted parts give a shortcut.  Since the first
$125$ vectors admit no closed set of $56$, any solution over all $137$
must select at least one of the $12$ vectors added by the last four
programs.  We split that condition into $12$ disjoint cases (vector $i$
selected, vectors $1,\dots,i-1$ not), each a smaller question with one
vector forced in.  All twelve are refuted.  Seven fell to a single kissat
run, in $27$ minutes to $2$ hours $51$ minutes; the later cases, which
exclude more of the new vectors, were the easiest.  Four needed a further
split into two or four parts on the most frequent free selection
variables, each part refuted in at most $3$ hours $32$ minutes.  The
first case, which forces one vector and excludes none, was refuted by an
eight-way split whose parts took $47$ minutes to $1$ hour $55$ minutes,
while one half of a two-way split took almost six hours.  The whole question, run directly under
both solvers and as four cubes, stayed undecided after four hours.  So
the record's $\Pt$ cannot reach $56$ from these programs either: at both
points, no recombination of the programs found lowers any side.

\begin{table}[ht]
\centering
\caption{Recombination questions.  \emph{Vectors} counts the dictionary
including the $16$ inputs; $c$ is the operation budget, one below the
side's best program.  Times are wall-clock on a machine running other
solver jobs in parallel.}\label{tab:verdicts}
\small
\begin{tabular}{llrrll}
\toprule
point & side & vectors & $c$ & verdict & how \\
\midrule
record & $L$ & $82$ & $54$ & no & MIP $90$\,s; SAT $0.4$\,s \\
record & $R$ & $81$ & $56$ & no & MIP $51$\,s \\
record & $R$ & $281$ & $56$ & no & $+200$ proposals; CaDiCaL ${\approx}30$\,s \\
record & $R$ & $2862$ & $56$ & open & structured; kissat $8$\,h, $\ell=24$ \\
balanced & $L$ & $86$ & $55$ & no & CaDiCaL ${\approx}30$\,s; MIP open at $6$\,h \\
balanced & $R$ & $92$ & $55$ & no & CaDiCaL ${\approx}30$\,s; MIP open at $6$\,h \\
balanced & $\Pt$ & $112$ & $56$ & no & order $306$\,s; levels $88$\,min (kissat) \\
record & $\Pt$ & $97$ & $56$ & no & order $33$\,s; levels $98$\,s \\
record & $\Pt$ & $111$ & $56$ & no & order $406$\,s \\
record & $\Pt$ & $125$ & $56$ & no & order $5000$\,s \\
record & $\Pt$ & $137$ & $56$ & no & $125$-vector part, then $12$ cases \\
record & $\Pt$ & $337$ & $56$ & open & $+200$ proposals; $8$\,h, levels \\
record & $\Pt$ & $3295$ & $56$ & open & structured; $8$\,h, $\ell=24$ \\
\bottomrule
\end{tabular}
\end{table}

\subsection{What a no means, and what comes next}

A refutation is a statement about a dictionary: no program built only
from those vectors is cheaper.  A $200$ could still exist that computes a
vector no search has produced, and these verdicts are unchecked solver
output so far.  The next steps are:
\begin{enumerate}
\item enlarge the dictionaries by proposals and structured vectors,
  starting from the open rows of Table~\ref{tab:verdicts}, with the same
  incremental split: each refuted dictionary forces any solution of a
  larger one to use a new vector;
\item certify every refutation used in a claim with a DRAT proof checked
  by an independent checker~\cite{DRATtrim};
\item try cutting-planes pseudo-Boolean solving~\cite{RoundingSat}, which
  reasons about the budget $\sum x_v\le c$ directly instead of through a
  counter's clauses; the converter to OPB is written;
\item seek lower bounds that apply to every program, not only to a
  dictionary: the floor of the record point is $156$, and a closure-level
  bound after reduction modulo $3$, where halvings become free sign
  changes, raises it only to $162$ (for programs whose scalings are
  powers of $2$), far below $201$;
\item look beyond this class.  Every operation count in
  Table~\ref{tab:rank48} belongs to the class descending from
  AlphaEvolve's scheme, and a second rank-48 class over $\Q$ would be a
  new starting point.
\end{enumerate}

\section{Data and code}

Everything is in the public repository
\url{https://github.com/gsidebottom/logic}; paths below are relative to
it.
\begin{itemize}
\item The scheme: \path{matmul/dps48/cand48/p005_201n_{L,R,P}.{slp,sms}},
  with \path{p005_201_Pt.slp}, the $57$-operation program for $\Pt$.  The
  same files are in PLinOpt as \path{data/4x4x4_48_201_{L,R,P}.{slp,sms}}.
\item Orbit annealing: \path{matmul/dps48/orbit48.py}.
\item Value-set search, recombination, and the SAT and MIP encodings:
  \path{src/bin/lns48.rs}.
\item Sweep and restarts: \path{matmul/dps48/patel/sweep.py} and
  \path{restarts.sh}.
\item Sign normalization and the exact verifier:
  \path{matmul/dps48/signnorm48.py} and \path{slp48check.py}.
\item Case splits and OPB conversion: \path{matmul/dps48/cube.py} and
  \path{cnf2opb.py}.
\item The Lean proofs: \path{proof/Mm/Correct48_201.lean} and
  \path{proof/Mm/Ring_4x4_48_201.lean}.  One of them runs alone with
\begin{verbatim}
cd proof && lake build Mm.Ring_4x4_48_201
\end{verbatim}
\item Appendix~\ref{app:scheme}: generated by
  \path{matmul/dps48/mk_paper_appendix48.py}, which re-checks the Brent
  equations first.
\end{itemize}

\section*{Acknowledgments}

This work was carried out in an extended interactive collaboration
with Claude (Anthropic), which implemented the search and verification
tools and much of this text under the author's direction and review.
Every claim is checked mechanically by the tools above; the results
rest on those checks rather than on trust in either the author or the
tools.  The method that found the scheme is Ruchit Patel's, whose
search code is public domain.

\begin{thebibliography}{99}
\bibitem{CaDiCaL} A.~Biere, T.~Faller, K.~Fazekas, M.~Fleury,
N.~Froleyks, F.~Pollitt, \emph{CaDiCaL 2.0}, Proc.\ CAV 2024, LNCS
14681, 133--152.
\bibitem{Kissat} A.~Biere, M.~Fleury et al., \emph{Kissat SAT solver},
\url{https://github.com/arminbiere/kissat} (version 4.0.4).
\bibitem{BCS97} P.~B\"urgisser, M.~Clausen, M.~A.~Shokrollahi,
\emph{Algebraic Complexity Theory}, Springer, 1997.
\bibitem{deGroote78} H.~F.~de~Groote, \emph{On varieties of optimal
algorithms for the computation of bilinear mappings I. The isotropy
group of a bilinear mapping}, Theoret.\ Comput.\ Sci.\ 7 (1978) 1--24.
\bibitem{Lean4} L.~de~Moura, S.~Ullrich, \emph{The Lean~4 theorem prover
and programming language}, Proc.\ CADE-28 (2021), LNCS 12699, 625--635.
\bibitem{DPS25} J.-G.~Dumas, C.~Pernet, A.~Sedoglavic, \emph{A
non-commutative algorithm for multiplying $4\times4$ matrices using
$48$ non-complex multiplications}, arXiv:2506.13242 (2025).
\bibitem{DPS26a} J.-G.~Dumas, C.~Pernet, A.~Sedoglavic, \emph{A more
accurate rational non-commutative algorithm for multiplying $4\times4$
matrices using $48$ multiplications}, arXiv:2603.18699 (2026).
\bibitem{DPSST26} J.-G.~Dumas, C.~Pernet, A.~Sedoglavic, J.~Stapleton,
P.~Tichavsk\'y, \emph{Fast matrix multiplication via recursive
$\langle4\times4\times4:48\rangle$ algorithms into practice},
arXiv:2609.12027 (2026).
\bibitem{PLinOpt} J.-G.~Dumas, B.~Grenet, C.~Pernet, A.~Sedoglavic,
\emph{PLinOpt, a collection of C++ routines handling linear \&
bilinear programs}, \url{https://github.com/jgdumas/plinopt} (data
files \texttt{4x4x4\_48\_202} and \texttt{4x4x4\_48\_201}, October 2026).
\bibitem{RoundingSat} J.~Elffers, J.~Nordstr\"om, \emph{Divide and
conquer: towards faster pseudo-Boolean solving}, Proc.\ IJCAI 2018,
1291--1299.
\bibitem{Fawzi22} A.~Fawzi et al., \emph{Discovering faster matrix
multiplication algorithms with reinforcement learning}, Nature 610
(2022) 47--53.
\bibitem{HiGHS} Q.~Huangfu, J.~A.~J.~Hall, \emph{Parallelizing the dual
revised simplex method}, Math.\ Program.\ Comput.\ 10 (2018) 119--142.
\bibitem{JT16} K.~Jamieson, A.~Talwalkar, \emph{Non-stochastic best arm
identification and hyperparameter optimization}, Proc.\ AISTATS 2016,
240--248.
\bibitem{Kaporin24} I.~E.~Kaporin, \emph{Finding complex-valued
solutions of Brent equations using nonlinear least squares}, Comput.\
Math.\ Math.\ Phys.\ 64 (2024) 1881--1891.
\bibitem{KM23} M.~Kauers, J.~Moosbauer, \emph{Flip graphs for matrix
multiplication}, Proc.\ ISSAC 2023; arXiv:2212.01175.
\bibitem{Mathlib} The mathlib Community, \emph{The Lean mathematical
library}, Proc.\ CPP 2020, 367--381.
\bibitem{MSY26} Y.~Moran, O.~Schwartz, S.~Yuan, \emph{Complex to
rational fast matrix multiplication}, arXiv:2602.13171 (2026).
\bibitem{Novikov25} A.~Novikov et al., \emph{AlphaEvolve: a coding agent
for scientific and algorithmic discovery}, arXiv:2506.13131 (2025).
\bibitem{Patel26} R.~Patel, \emph{matmul-4x4-48-202}, search code and
programs, \url{https://github.com/rp440/matmul-4x4-48-202} (October
2026, CC0).
\bibitem{Shaw98} P.~Shaw, \emph{Using constraint programming and local
search methods to solve vehicle routing problems}, Proc.\ CP 1998, LNCS
1520, 417--431.
\bibitem{Sinz05} C.~Sinz, \emph{Towards an optimal CNF encoding of
Boolean cardinality constraints}, Proc.\ CP 2005, LNCS 3709, 827--831.
\bibitem{Strassen69} V.~Strassen, \emph{Gaussian elimination is not
optimal}, Numer.\ Math.\ 13 (1969) 354--356.
\bibitem{DRATtrim} N.~Wetzler, M.~J.~H.~Heule, W.~A.~Hunt Jr.,
\emph{DRAT-trim: efficient checking and trimming using expressive
clausal proofs}, Proc.\ SAT 2014, LNCS 8561, 422--429.
\end{thebibliography}

\appendix
\section{The 201-operation programs}\label{app:scheme}
% generated by matmul/dps48/mk_paper_appendix48.py -- do not edit
The three programs below are the sign-normalized triple stored as
\path{matmul/dps48/cand48/p005_201n_{L,R,P}.slp} and in PLinOpt as
\path{data/4x4x4_48_201_{L,R,P}.slp}, in PLinOpt syntax.  In $L$
and $R$ the input \texttt{i}$(4i{+}j)$ is $a_{ij}$ (resp.\ $b_{ij}$), indices
from $0$, and output \texttt{o}$k$ is the factor of product $m_k$
($k=0,\dots,47$); in $P$ the input \texttt{i}$k$ is $m_k$ and output
\texttt{o}$(4i{+}j)$ is $c_{ij}$.  Each line is one operation (a binary
addition or subtraction, or a halving \texttt{/2}) or a free copy
\texttt{o}$k$\texttt{:=}$x$.  The products are $m_k=\ell_k\,r_k$ with
$\ell_k$ the $k$-th output of $L$ (left factor) and $r_k$ that of $R$.
The generator re-evaluates the programs over $\mathbb{Q}$ and checks all
$4096$ Brent equations before writing this appendix.

\subsection*{Left factors $L$: 53 additions + 2 halvings = 55 operations}
\begin{multicols}{4}\scriptsize
\begin{verbatim}
t1:=i9-i10;
t2:=i12-i15;
t3:=i6-i5;
t4:=i7+t3;
t5:=i0-i1;
t6:=i2+i3;
t7:=i4-i5;
t8:=i11-i8;
t9:=i14+i15;
t10:=t9-i13;
t11:=i3-i7;
t12:=i9-i13;
t13:=i2-i10;
t14:=t12/2;
t15:=t14+t9;
t16:=t14-t7;
t17:=t14-i10;
t18:=t11/2;
t19:=t18-t8;
t20:=t13-i14;
t21:=t18-t6;
t22:=t13-i6;
t23:=t19-i8;
t24:=i9-t14;
t25:=i4+t16;
t26:=t15+t21;
t27:=t5+t24;
t28:=t4+t18;
t29:=t26-t28;
t30:=t22+t25;
t31:=t17-t26;
t32:=t30+t26;
t33:=i0-t27;
t34:=t31-t20;
t35:=t29-t25;
t36:=t32+t23;
t37:=t31+t23;
t38:=t29+t33;
t39:=t32+t34;
t40:=i12+t34;
t41:=t30-t33;
t42:=t16+t40;
t43:=t27-t42;
t44:=t43+t34;
t45:=t19-t42;
t46:=t43-t40;
t47:=t45+t25;
t48:=t46-t45;
t49:=t14-t45;
t50:=t2+t46;
t51:=t48+t42;
t52:=t48+t33;
t53:=t18-t51;
t54:=t48+t24;
t55:=t51+t23;
o0:=t1;
o1:=t2;
o2:=t4;
o3:=t51;
o4:=t35;
o5:=t52;
o6:=t36;
o7:=i4;
o8:=t15;
o9:=i2;
o10:=t47;
o11:=t27;
o12:=t6;
o13:=t53;
o14:=t28;
o15:=i0;
o16:=t43;
o17:=t54;
o18:=t48;
o19:=t49;
o20:=t8;
o21:=t16;
o22:=t37;
o23:=t9;
o24:=t29;
o25:=i10;
o26:=t44;
o27:=t17;
o28:=t10;
o29:=t55;
o30:=i12;
o31:=t45;
o32:=t38;
o33:=t19;
o34:=t30;
o35:=t7;
o36:=t20;
o37:=t31;
o38:=t39;
o39:=t40;
o40:=t5;
o41:=i8;
o42:=t3;
o43:=t21;
o44:=t41;
o45:=t22;
o46:=t50;
o47:=t32;
\end{verbatim}
\end{multicols}

\subsection*{Right factors $R$: 55 additions + 2 halvings = 57 operations}
\begin{multicols}{4}\scriptsize
\begin{verbatim}
t1:=i1+i5;
t2:=i9+t1;
t3:=i0+i4;
t4:=i14-i10;
t5:=i2+i6;
t6:=i0+i12;
t7:=i3+i15;
t8:=i3+i7;
t9:=t6-i8;
t10:=i11+t8;
t11:=i9-i13;
t12:=i1-i3;
t13:=i8+i10;
t14:=i14+i15;
t15:=t14-i13;
t16:=t14-i12;
t17:=t11-i1;
t18:=t12/2;
t19:=t13/2;
t20:=i8-t19;
t21:=t7+t18;
t22:=i6+t18;
t23:=t19-i9;
t24:=t15-t23;
t25:=t17+t18;
t26:=t23-t1;
t27:=t10-t20;
t28:=t6-t19;
t29:=i14-t19;
t30:=t22-t5;
t31:=t8-t27;
t32:=t28-t21;
t33:=t29-t32;
t34:=t24+t32;
t35:=t30-t33;
t36:=t24+t31;
t37:=t25+t32;
t38:=t23+t37;
t39:=t34-t30;
t40:=t16-t33;
t41:=t37-t31;
t42:=t34-t40;
t43:=t3-t40;
t44:=t43-t27;
t45:=t44-t22;
t46:=t26+t44;
t47:=t43+t46;
t48:=t45-t31;
t49:=t47-t45;
t50:=t46+t40;
t51:=t47-i4;
t52:=t49-t44;
t53:=t49+t19;
t54:=t45+t20;
t55:=t52+t18;
t56:=t52-t30;
t57:=t49-t23;
o0:=t48;
o1:=t42;
o2:=t15;
o3:=t21;
o4:=t2;
o5:=t3;
o6:=t38;
o7:=t56;
o8:=t49;
o9:=t4;
o10:=t50;
o11:=t22;
o12:=i14;
o13:=i15;
o14:=t24;
o15:=t5;
o16:=t25;
o17:=i4;
o18:=t43;
o19:=t51;
o20:=t16;
o21:=t52;
o22:=t6;
o23:=t53;
o24:=t26;
o25:=t54;
o26:=t17;
o27:=t45;
o28:=t57;
o29:=t7;
o30:=t39;
o31:=t46;
o32:=t1;
o33:=t33;
o34:=t27;
o35:=t55;
o36:=t9;
o37:=t28;
o38:=t41;
o39:=t34;
o40:=i6;
o41:=t35;
o42:=t36;
o43:=t29;
o44:=t8;
o45:=t10;
o46:=t11;
o47:=t37;
\end{verbatim}
\end{multicols}

\subsection*{Output map $P$: 89 additions = 89 operations}
\begin{multicols}{4}\scriptsize
\begin{verbatim}
t1:=i43+i12;
t2:=i9-i12;
t3:=t2+i40;
t4:=i11-t3;
t5:=i15-t3;
t6:=i3+i13;
t7:=i29+i13;
t8:=t4+i31;
t9:=t8+i19;
t10:=i21+i35;
t11:=i27+t6;
t12:=t11+i25;
t13:=t2-t12;
t14:=t9+t13;
t15:=t10-t14;
t16:=t12+t15;
t17:=i17+t16;
t18:=t17-i18;
t19:=i20-t18;
t20:=i25+i0;
t21:=t18-t20;
t22:=i33+t19;
t23:=t1-t15;
t24:=i8-t14;
t25:=t24-i23;
t26:=i16+t25;
t27:=i34+t22;
t28:=t25-t22;
t29:=t28+t23;
t30:=t23-i47;
t31:=t13+t29;
t32:=i46-t26;
t33:=i26-t26;
t34:=t18+t32;
t35:=t6+t1;
t36:=t35+t31;
t37:=i37+t4;
t38:=t37+t31;
t39:=i39+t36;
t40:=i24+t36;
t41:=i28-i23;
t42:=t41+t34;
t43:=t34-i1;
t44:=t43+t39;
t45:=t27-i45;
t46:=i4+t40;
t47:=t44-t17;
t48:=t47+i5;
t49:=i10-t8;
t50:=t49+t45;
t51:=t50+t44;
t52:=t51-t46;
t53:=i14-t4;
t54:=t53+t29;
t55:=t54+t32;
t56:=i21+i7;
t57:=t56-t46;
t58:=t57-t45;
t59:=t55-i42;
t60:=t10-i2;
t61:=t60+t55;
t62:=t59+t45;
t63:=i36-t16;
t64:=t63-t38;
t65:=t38-i22;
t66:=t44+t64;
t67:=t33-t64;
t68:=t42+t33;
t69:=t65+t7;
t70:=t42-t61;
t71:=t61-t46;
t72:=t59+t21;
t73:=t21-t7;
t74:=i32+t40;
t75:=t74+t32;
t76:=t75+t70;
t77:=i30-t39;
t78:=t77+t67;
t79:=t27-i44;
t80:=t79+t72;
t81:=t19-i41;
t82:=t81+t69;
t83:=i38-t67;
t84:=t83+t72;
t85:=t84+t30;
t86:=i6-t9;
t87:=t86+t70;
t88:=t87+t69;
t89:=t88+t30;
o0:=t48;
o1:=t76;
o2:=t5;
o3:=t80;
o4:=t52;
o5:=t71;
o6:=t58;
o7:=t62;
o8:=t65;
o9:=t89;
o10:=t82;
o11:=t73;
o12:=t66;
o13:=t68;
o14:=t78;
o15:=t85;
\end{verbatim}
\end{multicols}


\end{document}
